\section{Introduction}

The paradigm of Large Language Model (LLM) training has undergone a pivotal shift from Supervised Fine-Tuning (SFT) to Reinforcement Learning-based (RL) alignment. While models such as GPT-4 \cite{Achiam2023GPT4TR} and Claude \cite{TheC3} have demonstrated outstanding reasoning capabilities, their reliance on complex reward models imposes significant computational and scalability bottlenecks. Recently, the release of DeepSeek-R1 popularized Group Relative Policy Optimization (GRPO) \cite{deepseekmath}, marking a milestone in training specialized reasoning models without the heavy overhead of a traditional reward model. By leveraging group-relative rewards, GRPO allows the policy model to autonomously distinguish between good and bad reasoning paths within the same prompt context, effectively optimizing for complex system thinking capabilities.

However, as models scale and tasks become more complex, such as Olympiad-level mathematics and complex reasoning chains, vanilla GRPO faces diminishing returns. We identify three fundamental limitations in existing RL pipelines for reasoning:

\textbf{1. Zero Advantage:} In many reasoning tasks, especially at the start or late stages of training, a model may get all rollouts in a group share the same outcome, either all correct or all incorrect. In such cases, the mean reward equals the individual reward, and the standard deviation approaches zero, causing the advantage signal become zero \cite{zhang2026trainlesslearnmore}. This stalls the gradient and wastes expensive compute.

\textbf{2. The Compulsive Correction Fallacy:} Multi-layer \cite{mgrpo} or iterative reflection models \cite{madaan2023selfrefineiterativerefinementselffeedback} often suffer from overcorrection. When prompted to reflect, models exhibit a systematic bias towards altering their initial answers, even when correct. This bias degrades the stability of the model outputs and leads to a high rate of false-negative revisions.

\textbf{3. The Noise-Signal Dichotomy:} Traditional RL update paradigms implicitly operate under the assumption that all optimization errors possess equivalent informational value. However, a stochastic error originating from a superficial computational oversight yields significantly less pedagogical utility than a structural error rooted in a fundamental logical misconception. Standard GRPO frameworks currently lack the granular mechanism required to disambiguate these distinct failure modes.

To address these limitations, we propose Semantic-Balanced GRPO (SB-GRPO). Our framework tackles the zero advantaged problem by introducing a Heuristic Confidence Gate, which applies a differentiable switching mechanism to activate Negative-enhanced GRPO (NGRPO) \cite{ngrpo} using virtual rewards for uniform batches. To mitigate correction bias and differentiate noise from systematic errors, we integrate a Semantic Error Profiling (SEP) module within a two-layer generation-reflection pipeline. By clustering reasoning paths and constructing Static Balanced Batches, our approach ensures that the model optimizes based on the most highly informative errors rather than trivial computational noise.

% \section{Related Work}
% Reinforcement Learning from Human Feedback (RLHF) \cite{10.5555/3294996.3295184}\cite{Ji2024ReinforcementLF} has predominantly relied on Proximal Policy Optimization (PPO) \cite{ppo} to align LLMs. While computationally robust, PPO inherently relies on dedicated reward models to guide the policy optimization process, thereby severely increasing the memory footprint and computational overhead during training. To circumvent this bottleneck, GRPO \cite{deepseekmath} was introduced, which effectively eliminates the parameter overhead of the critic by substituting it with a baseline derived from the average rewards of $G$ distinct rollouts. Building upon this memory-efficient paradigm, Negative-enhanced GRPO (NGRPO) \cite{ngrpo} further refined the objective function by incorporating virtual negative rewards. This modification ensures continuous gradient flow and mitigates optimization stalling in extreme edge cases where no correct reasoning trajectories are generated within a sampled group.

% In parallel with advancements in optimization efficiency, augmenting the intrinsic reasoning capabilities of LLMs via explicit reasoning paths such as Chain-of-Thought \cite{10.5555/3600270.3602070} and autonomous self-correction mechanisms \cite{madaan2023selfrefineiterativerefinementselffeedback} has garnered significant research attention. Within this context, Multi-Layer GRPO (MGRPO) \cite{mgrpo} demonstrated that such self-correction can be effectively internalized through a two-layer generation pipeline. By utilizing the generated output from the first layer as the input context for the second layer, models can effectively verify and correct their own errors. However, a critical vulnerability of MGRPO and similar unconstrained self-reflective frameworks is their tendency to induce a pronounced correction bias. This phenomenon coerces models to arbitrarily alter initially valid reasoning paths, consequently degrading inference stability and limiting the overall efficacy of the self-correction loop.

% To address the instability and overcorrection inherent in stochastic reasoning frameworks, recent studies have increasingly focused on uncertainty quantification and sophisticated data curation techniques. Semantic Entropy (SE) \cite{kuhn2023semantic} provides a robust information-theoretic metric for distinguishing genuine model confusion from generative diversity. This concept was leveraged by SEED-GRPO \cite{seedgrpo}, which integrated SE into the loss function to dynamically dampen gradient updates for highly uncertain prompts. Concurrently, theoretical insights from Balanced Sampling \cite{divagrpo} emphasize that maintaining an equitable distribution of correct and hard-negative samples is pivotal for stabilizing reinforcement learning convergence in sparse-reward environments. Collectively, these isolated advancements underscore the necessity of a unified framework—such as our proposed approach—that synergizes semantic error profiling with statically balanced batches to robustly optimize iterative reasoning.